\section{Atomistic study of Positive Exchange Bias in Hard/Soft ferromagnetic bilayer nanosystem \href{https://doi.org/10.1088/1402-4896/acb23c}{[\textit{Physica Scripta, 98 2 025403 (2023)}]}}


\section*{Main points from the Paper}
\subsection*{\centering\underline{Abstract}}
\hspace{2cm} The calculations for the exchange spring nano-bilayer with intermixing of high and low anisotropic moments at the interface are studied in this paper, and an unexpected phenomenon of positive exchange bias is observed. \hl{The influence of temperature is clearly observed in the shift and the width of the M-H loop}. They have claimed that the interface roughness is the main cause for the shift in the hysteresis loop.    
\subsection*{\underline{I. Introduction}}
\begin{enumerate}
    \item The multilayer exhibits perpendicular magnetic anisotropy along with an out-of-plane exchange bias field. They have observed a monotonically increasing EB magnitude with the annealing field, which saturates at $\pm1$ T.
    \item The exchange-bias effect was seen first in the core-shell structure, and later, this was also observed for the FM-AFM system.
    \item The coupling between FM and AFM spins at the interface adds an internal field (bias field) to the applied field, which is a trivial source of asymmetrical hysteresis loop and eventually leads to the existence of exchange bias. 
    \item Also, the induction of anisotropy at the interface might be another reason for the unidirectional shift of the loop. [\textcolor{red}{ref. 4 -anisotropy affects the EB}]
    \item Exchange bias is commonly used in magnetic sensors due to the enhanced thermal stability of the material; other technological applications include GMR devices, hard disks, and hard drives, and it also serves as the basis for many other applications in spintronics.
    \item Of many existing theoretical assumptions for the exchange bias, one of the main ones is the presence of interfacial uncompensated AFM spin magnetic moments [\textcolor{red}{ref.9: Exchange Bias paper.}]
    \item Different morphology systems have been studied for EB-related studies, mostly done for the thin film bilayers, as there can be enhanced control over the interface, and they are more suitable for device fabrication.
    \item The shift in the M-H loop can be in any direction.
    \item Positive EB is far less common than negative EB as found in several heterostructures.
    \item The presence of AFM exchange coupling at the interface of the FM bilayer system is responsible for the EB effect. 
    \item A large amount of disorder present at the interface of the FM-AFM can also be the reason for this effect.
    \item Atomistic spin models are used to study the exchange-bias field caused by the roughness in simple cubic Co/CoO core-shell NPs \cite{wuInfluenceAtomicRoughness2020,influenceofinterfacial2011}. 
    \item A key feature of the exchange bias effect, i.e., unidirectional shift of the hysteresis loop, is observed.
\end{enumerate}

\subsection*{\underline{II. Methods}}
\begin{enumerate}
    \item Classical Heisenberg Hamiltonian and Landau-Lifshitz Gilbert equations are employed to study the EB effect in Hard/Soft FM Fe/Co bilayers.
    \item Both flat and roughened surfaces are simulated, particularly looking for the effect of roughness on the EB.
    \item Details about the system for the simulation and system are:
    \begin{enumerate} [\ding{212}]
        \item BCC crystal structure for both Fe and Co.
        \item Dimension of the system: $x=10$ nm, $\times y=10$ nm, $\times z=4$ nm. total atoms, 32368 atoms with $50\%$ Fe and $50\%$ Co.
        \item Understand the extensively written Hamiltonian.
        \item Exchange interactions: {$J_{Co}=-6.064\times10^{-21}J/link$}; \\{$J_{Fe}=-7.05\times10^{-21}J/link$}; {$J_{Fe-Co}=1.0\times10^{-22}J/link$}
        \item Uniaxial anisotropy: {$K_{Co}=6.69\times10^{-24}J/atom$}; \\{$K_{Fe}=5.65\times10^{-25}J/atom$}
        \begin{table}
            \centering
            \begin{tabular}{ccccc}\hline \hline
                Symbol & Value & Parameter & Units & \\ \hline 
                $\alpha$ &  & Gilbert Damping paramter  & - & \\
                 &  &  &  & \\
                 &  &  &  & \\
                 &  &  &  & \\
                 &  &  &  & \\
                 &  &  &  & \\ \hline \hline
            \end{tabular} 
            \caption{Parameters used for the Exchange bias calculation for the Fe/Co in Junaid Ahsan's paper}
            \label{tab:Fe/Co_parameter_EB}
        \end{table}
        \item Uniaxial anisotropy direction = 0, 0, 1 for Co and 1, 0, 0 for Fe.
        \item Heun integration scheme was used for the hysteresis calculations.
        \item Here for the easy damping, the critical damping value is used; $\alpha=1.0$ [\textcolor{red}{19, 20}]
        \item Applied field was varied from: $+1.0$ T to $-1.0$ T in steps of $0.001$ T.
        \item The magnetic field was applied along the $xy-plane$ of the bilayer.
        \item The details of time-steps are;
        \begin{verbatim}
            equilibration-time-steps=100000
            loop-time-steps=100000
        \end{verbatim}   
    \end{enumerate}
\end{enumerate}

\subsection*{\underline{III. Results and Discussions}}
\begin{enumerate}
    \item The influence of interface roughness was studied by Evans et.al.,\cite{influenceofinterfacial2011} for FM-AFM core-shell spherical nanoparticles. The roughness gives rise to uncompensated magnetic moments at the interface. 
    \item The degree of coupling between the FM and AFM was increased by the roughness at the interface. 
    \item EB fundamentally being an interfacial phenomenon, any modification of the interface structure has a significant effect on it.
    \item The change in the microstructure in the different crystalline orientations changes due to intermixing of hard and soft atomic spin moments.
    \item Additionally, the other prerequisite for EB is that there is a difference in the anisotropy of two catenated layers; if not, then they will rotate under the effect of an externally applied field coherently.
    \item In this paper, the shape of the M-H loop is decided by the direction of the applied field and the magnetic anisotropy in the Fe and Co.
    \item The uniaxial anisotropy direction is set as per the easy axes of the individual system. For Fe, it is 100; for Co, it is 001. This is the case for a bulk system.
    \item If the anisotropy is along the same direction, then one would expect a more square loop.
    \item There is a difference in the magnitude and also the direction of the anisotropy in this paper for observing the Positive EB (PEB).
    \item There is a shift in a positive direction in all of the cases when the interface is flat, i.e., there is no intermixing.
    \item The positive EB in the smooth interfaced bilayer can thus be attributed to the;
    \begin{itemize}
        \item larger magnetic anisotropy energy of the hard (Co) layer than the soft (Fe) layer.
        \item ferrimagnetic nature of the interface.
        \item anisotropy direction.
    \end{itemize}
    \item The alloy formation at the interface leads to the enhancement in the EB as observed in all simulated temperatures.
    \item There is a decrease in the coercivity with elevating temperature.
\end{enumerate}

\subsection*{\underline{IV. Conclusions and further work}}
\begin{enumerate}
    \item Positive EB was studied using ASD, with AFM coupled Fe/Co magnetic bilayer system.
    \item It is verified that EB is an interfacial phenomenon.
    \item Major contribution comes from the interfacial exchange interaction between Co and Fe magnetic spin moments.
    \item Improvements in the positive EB were observed after increasing the roughness.
\end{enumerate}
\vspace{1cm}
The effect of EB variation on the degree of roughness and interfacial mixing of atomic moments can be studied. - \hl{Future Work Related}